\section{Considerações finais}
\label{sec:consideracoes}

O pitch da \projeto{} foi estruturado em onze lâminas que cobrem os oito tópicos exigidos pela orientação e acrescentam abertura em \textit{storytelling}, equipe e pedido. A narrativa de Maria e Thor conduz o público do problema à solução; os dados são apresentados com fonte e as projeções, com o devido rótulo de hipótese; e a divisão da fala acompanha a estrutura organizacional do grupo.

A principal limitação a reconhecer é que a tração do negócio ainda é de natureza qualitativa: cinco sessões de teste guiado, um protótipo funcional e um modelo econômico baseado em estimativas. O pitch, por isso, apresenta o negócio como startup em fase de validação, e não como operação com receita. Essa postura é deliberada: uma apresentação que declara o que foi validado e o que ainda será testado tende a ser mais crível do que uma que sugere maturidade inexistente.
